% Part: incompleteness
% Chapter: theories-computability
% Section: oconsis-ext-of-q-undec

\documentclass[../../../include/open-logic-section]{subfiles}

\begin{document}

\olfileid{inc}{tcp}{oqn}

\olsection{$\Th{Q}$ चा $\omega$-सुसंगत विस्तार अनिर्णेय असतो}

\begin{explain}
$\Th{Q}$ हा c.e.-संपूर्ण आहे या सिद्धतेत,
$\Th{Q}$ मध्ये सिद्ध होणारे प्रत्येक वाक्य नैसर्गिक संख्यांबाबत
``सत्य'' असते या तथ्याचा आधार घेतला होता. पुढील व्याख्या आणि
प्रमेय, वरील सिद्धतेत आवश्यक असलेले ``सत्य''चे नेमके पैलू
ओळखून, त्या प्रमेयाचे सबलीकरण करतात. मात्र या प्रमेयावर फार
रेंगाळू नका, कारण लवकरच आपण त्याहून अधिक सबल प्रमेय सिद्ध करू.
ते मुख्यतः ऐतिहासिक कारणासाठी समाविष्ट केले आहे: गोडेल यांच्या
मूळ लेखात $\omega$-सुसंगततेची संकल्पना वापरली होती; पण लगेचच
$\omega$-सुसंगततेच्या जागी साधी सुसंगतता घालून त्यांचा निकाल
अधिक सबल करण्यात आला.
\end{explain}

\begin{defn}
\ollabel{thm:oconsis-q}
उपपत्ती $\Th{T}$ ही $\omega$-सुसंगत असते, जर पुढील अट पूर्ण होत असेल:
$\lexists[x][!A(x)]$ हे कोणतेही वाक्य असो. $\Th{T}$ कडून
$\lnot !A(\num 0)$, $\lnot !A(\num 1)$,
$\lnot !A(\num 2)$, \dots ही सर्व वाक्ये सिद्ध होत असतील,
तर $\Th{T}$ कडून $\lexists[x][!A(x)]$ सिद्ध होत नाही.
\end{defn}

\begin{thm}
\ollabel{thm:oconsis-q-undec}
  $\Th{T}$ ही कोणतीही $\omega$-सुसंगत उपपत्ती असेल आणि तिच्यात
  $\Th{Q}$ समाविष्ट असेल, तर $\Th{T}$ अनिर्णेय असते.
\end{thm}

\begin{proof}
$\Th{T}$ मध्ये $\Th{Q}$ समाविष्ट असल्यास संगणनक्षम
फलनांचे आणि संबंधांचे $\Th{T}$ मध्ये निरूपण होते.
आधीची सिद्धता किंचित बदलली तरी पुरेसे आहे. पूर्वीप्रमाणे,
$x \in K$ असेल, तर $\Th{T}$ हे $\lexists[s][!A_T(\num
  x,\num x, s)]$ सिद्ध करते. उलट, $\Th{T}$ हे
$\lexists[s][!A_T(\num x, \num x, s)]$ सिद्ध करते असे धरू.
मग $x$ हा $K$ मध्ये असलाच पाहिजे: अन्यथा यंत्र $x$ चे
$x$ या आदानावरील थांबणारे संगणन अस्तित्वात नसते. $!A_T$ हे
क्लिनीच्या $T$ संबंधाचे निरूपण करत असल्याने $\Th{T}$ कडून
$\lnot !A_T(\num x, \num x, \num 0)$,
$\lnot !A_T(\num x, \num x, \num 1)$, \dots
सिद्ध होतील. यामुळे $\Th{T}$ ही $\omega$-विसंगत ठरेल.
\end{proof}

\end{document}
